\documentclass[11pt,a4paper,fontset=fandol]{ctexart}

\usepackage[a4paper,top=2.4cm,bottom=2.6cm,left=2.35cm,right=2.35cm,headheight=18pt,headsep=0.55cm,footskip=26pt]{geometry}
\usepackage{amsmath,amssymb}
\usepackage{fancyhdr}
\usepackage{lastpage}
\usepackage{enumitem}
\usepackage{titlesec}
\usepackage{xcolor}
\usepackage{hyperref}
\usepackage{bookmark}
\usepackage[most]{tcolorbox}
\usepackage{setspace}
\usepackage{array}

\hypersetup{colorlinks=true,linkcolor=blue!50!black,urlcolor=blue!50!black}

\setstretch{1.16}
\setlength{\parindent}{2em}
\setlength{\parskip}{0.32em}
\setlist[enumerate]{itemsep=0.35em, topsep=0.35em, leftmargin=2.3em}
\setlist[itemize]{itemsep=0.25em, topsep=0.25em, leftmargin=2em}

\definecolor{mainblue}{RGB}{36,83,140}
\definecolor{lightblue}{RGB}{241,246,252}
\definecolor{lightgray}{RGB}{246,246,246}
\definecolor{rulegray}{RGB}{110,118,128}
\definecolor{softgreen}{RGB}{236,247,240}
\definecolor{softyellow}{RGB}{254,249,229}

\titleformat{\section}{\Large\bfseries\color{mainblue}}{\thesection.}{0.45em}{}[\vspace{0.25em}\color{rulegray}\titlerule]
\titlespacing*{\section}{0pt}{1.2em}{0.8em}
\titleformat{\subsection}{\large\bfseries\color{mainblue}}{\thesubsection}{0.5em}{}

\pagestyle{fancy}
\fancyhf{}
\fancyhead[L]{\small 欧拉路径与哈密顿路径专题目录页}
\fancyhead[C]{\small 组合专题资料索引}
\fancyhead[R]{\small 刘文博}
\fancyfoot[C]{\small 第 \thepage 页 / 共 \pageref{LastPage} 页}
\renewcommand{\headrulewidth}{0.55pt}
\renewcommand{\footrulewidth}{0.35pt}

\newtcolorbox{goalbox}[1]{breakable,colback=lightblue,colframe=mainblue,boxrule=0.7pt,arc=1mm,title=\textbf{#1},fonttitle=\bfseries}
\newtcolorbox{infobox}[1]{breakable,colback=softgreen,colframe=green!40!black,boxrule=0.65pt,arc=1mm,title=\textbf{#1},fonttitle=\bfseries}
\newtcolorbox{warnbox}[1]{breakable,colback=softyellow,colframe=orange!60!black,boxrule=0.65pt,arc=1mm,title=\textbf{#1},fonttitle=\bfseries}

\begin{document}

\thispagestyle{empty}
\begin{titlepage}
\centering
\vspace*{0.9cm}

\begin{tcolorbox}[enhanced,width=0.92\textwidth,colback=white,colframe=mainblue,boxrule=1pt,arc=0mm,left=6mm,right=6mm,top=8mm,bottom=8mm]
\centering

{\fontsize{24}{30}\selectfont\bfseries 欧拉路径与哈密顿路径专题目录页\par}
\vspace{0.65cm}
{\fontsize{17}{22}\selectfont 讲义、题单与周测的使用说明\par}

\vspace{1cm}
\begin{tcolorbox}[colback=lightblue,colframe=mainblue,width=0.84\textwidth,boxrule=1pt]
\centering
{\large 适合对象：已经学过基础图论与染色方法，准备系统提升“图上走法”判断能力的学生}\\[0.35em]
{\normalsize 本目录页帮助把欧拉与哈密顿专题资料串成一个完整训练包}
\end{tcolorbox}

\vspace{1.0cm}
\renewcommand{\arraystretch}{1.42}
\begin{tabular}{>{\bfseries}p{3.5cm}p{8cm}}
专题名称： & 欧拉路径与哈密顿路径 \\
资料组成： & 课堂讲义、提高训练题单、周测 A 卷、周测 B 卷 \\
使用方式： & 先学讲义，再做题单，最后用 A/B 周测检查掌握情况 \\
编译方式： & XeLaTeX \\
\end{tabular}

\vspace{0.9cm}
{\large \today\par}
\end{tcolorbox}
\end{titlepage}

\tableofcontents
\newpage

\section{专题包包含哪些资料}

\begin{goalbox}{四份核心资料}
\begin{enumerate}
  \item 课堂讲义：帮助分清欧拉与哈密顿两类问题，建立最核心的判定框架；
  \item 提高训练题单：用于课后巩固和变式提升，重点训练“会区分模型、会排除、会构造”；
  \item 周测 A 卷：检查欧拉判定与基础哈密顿必要条件是否稳定；
  \item 周测 B 卷：检查结构分析、反例构造和综合判断是否真正到位。
\end{enumerate}
\end{goalbox}

\begin{infobox}{对应文件名}
\begin{itemize}
  \item 讲义：\texttt{euler\_hamilton\_handout.tex} / \texttt{euler\_hamilton\_handout.pdf}
  \item 提高训练题单：\texttt{euler\_hamilton\_advanced\_problem\_set.tex} / \texttt{euler\_hamilton\_advanced\_problem\_set.pdf}
  \item 周测 A 卷：\texttt{euler\_hamilton\_weekly\_test\_A.tex} / \texttt{euler\_hamilton\_weekly\_test\_A.pdf}
  \item 周测 B 卷：\texttt{euler\_hamilton\_weekly\_test\_B.tex} / \texttt{euler\_hamilton\_weekly\_test\_B.pdf}
\end{itemize}
\end{infobox}

\section{建议学习顺序}

\begin{goalbox}{推荐顺序}
\begin{enumerate}
  \item 先完整学习讲义，重点吃透“欧拉看边、哈密顿看点”这个最核心区别；
  \item 再完成提高训练题单，把奇度点判定、二分图染色和结构排除做熟；
  \item 接着做周测 A 卷，检查基础判断是否稳定；
  \item 最后做周测 B 卷，检查能否做综合说明和反例构造。
\end{enumerate}
\end{goalbox}

\begin{warnbox}{不建议的做法}
\begin{itemize}
  \item 看到“走法题”就只盯着度数，不先分清题目到底看边还是看点；
  \item 把哈密顿问题也机械套成“奇度点判定”；
  \item 满足必要条件后就直接断言“一定存在”。
\end{itemize}
\end{warnbox}

\section{每份资料主要解决什么问题}

\subsection{课堂讲义}
讲义适合第一次系统学习这个专题，主要解决以下问题：
\begin{itemize}
  \item 欧拉道路、欧拉回路、哈密顿路径、哈密顿回路分别是什么意思；
  \item 为什么欧拉问题由度数奇偶决定；
  \item 为什么哈密顿问题没有同样简单的统一判定法；
  \item 怎样用二分图染色和结构切分排除哈密顿回路。
\end{itemize}

\subsection{提高训练题单}
题单适合第二阶段训练，主要用来把方法做熟。它更强调：
\begin{itemize}
  \item 能否快速判断题目属于欧拉还是哈密顿；
  \item 能否正确使用奇度点结论；
  \item 能否用度数为 $1$ 的点、二分图黑白点数、公共顶点等结构进行排除；
  \item 能否举出合适反例说明“必要条件不充分”。
\end{itemize}

\subsection{周测 A 卷}
A 卷偏基础，适合检查：
\begin{itemize}
  \item 欧拉道路与欧拉回路的判定是否稳定；
  \item 基础哈密顿路径与回路是否会判断；
  \item 是否能用最经典的黑白染色处理网格图。
\end{itemize}

\subsection{周测 B 卷}
B 卷偏综合，适合检查：
\begin{itemize}
  \item 是否会举反例区分两类问题；
  \item 是否会进行结构说明，而不只是写结论；
  \item 是否真正理解必要条件与充分条件的差别。
\end{itemize}

\section{一个推荐的 4 次训练安排}

\begin{goalbox}{4 次完成一个小专题}
\begin{enumerate}
  \item 第 1 次：学习讲义前半部分，稳住欧拉判定；
  \item 第 2 次：学习讲义后半部分，重点理解哈密顿问题的常见排除方法；
  \item 第 3 次：完成提高训练题单，并把错题按“模型分错 / 判定错 / 构造不会 / 证明不完整”分类订正；
  \item 第 4 次：完成周测 A 卷或 B 卷，作为阶段性检查。
\end{enumerate}
\end{goalbox}

\begin{infobox}{如果孩子基础较好，可以这样压缩}
\begin{itemize}
  \item 第 1 天：讲义 + 题单基础题；
  \item 第 2 天：题单综合题 + 错题订正；
  \item 第 3 天：周测 A 卷；
  \item 第 4 天：周测 B 卷。
\end{itemize}
\end{infobox}

\section{专题学习后的自检清单}

\begin{goalbox}{如果下面 8 条都能做到，说明这个专题基本过关}
\begin{enumerate}
  \item 我能准确说出欧拉问题和哈密顿问题的区别；
  \item 我会判断连通图是否有欧拉回路；
  \item 我会判断连通图是否有欧拉道路；
  \item 我知道哈密顿问题不能机械地套奇度点结论；
  \item 我会用黑白染色排除某些网格图的哈密顿回路；
  \item 我会利用度数为 $1$ 的点排除哈密顿路径；
  \item 我会举出“有欧拉无哈密顿”或“有哈密顿无欧拉”的例子；
  \item 我知道必要条件满足后，往往还需要具体构造或进一步结构分析。
\end{enumerate}
\end{goalbox}

\section{给家长的使用建议}

\begin{warnbox}{更有效的带练方式}
\begin{itemize}
  \item 每做一道题，先让孩子先说“这题看边还是看点”；
  \item 如果卡住，可以只提示“数数奇度点”或“试试黑白染色”；
  \item 订正时重点不只是对错，更要看孩子是否把两类模型混淆；
  \item 反例题建议让孩子先自己画图，再比较不同例子。
\end{itemize}
\end{warnbox}

\begin{infobox}{和后续专题的衔接}
这个专题学完以后，可以继续自然过渡到二分图、匹配与覆盖、更系统的网格图路径问题，以及图论中的构造与极值专题。
\end{infobox}

\end{document}
